\section{Worked Cases I: Calendar and Precedence}
\label{sec:calendar-cases}

These cases are stated as conditions rather than assigned to an unchecked civil year.  Each answer begins with the calendar, not with the desired formulary.  A “proper” feast below exists only when the approved calendar named in the case actually contains it at the stated rank.

\begin{casefile}{1}{The Immaculate Conception on an Advent Sunday}
\item[Calendar facts] December 8 is both the universal I-class Immaculate Conception and a I-class Sunday of Advent.
\item[Ranks and rows] The feast is position 4; the Sunday is position 6.  Rubric 15 also states the exception expressly.
\item[Impeded item] The Advent Sunday is privilegedly commemorated; it is not anticipated, resumed, or transferred.
\item[Mass category] This is the Mass of the occurring Office, not a Marian votive.
\item[Assembly] Use the complete proper Mass of the Immaculate Conception in white, with Gloria and Creed.  Add the Sunday Collect, Secret, and Postcommunion as the one privileged commemoration, in that order.  The feast's Preface rule controls; the commemoration does not import a Sunday Preface or Gospel.  Use the ordinary Canon and ordinary Last Gospel unless the Missal assigns a proper intervention.
\item[Boundary] This exceptional result belongs to the Immaculate Conception; it does not create a general rule that a Marian feast defeats Advent Sunday.
\item[Rubric locus] 15; precedence table 91 positions 4 and 6; 109a, 111a, 113; 431--432, 475--483, 505, 509.
\end{casefile}

\begin{casefile}{2}{The Annunciation impeded through the Paschal sequence}
\item[Calendar facts] March 25 falls on Palm Sunday, or so near Easter that the I-class Annunciation cannot be placed before Easter without meeting superior days.
\item[Ranks and rows] The Annunciation is a universal I-class feast at position 11; Palm Sunday is position 6, Holy Week ferias position 7, the Triduum position 2, Easter position 1, Easter-octave days position 10, and Low Sunday position 6.
\item[Impeded item] The Annunciation transfers.  When its transfer must pass Easter, Monday after Low Sunday is its named proper seat.  It is not also commemorated on the impeding day.
\item[Mass category] Celebrate each impeding temporal day from its own Office.  On the named Monday celebrate the transferred feast as a I-class Mass of the Office.
\item[Assembly] Palm Sunday and Holy Week retain their unique temporal formularies and rites.  The transferred Annunciation uses its complete white proper, Gloria, Creed, assigned Preface, and its normal Canon and conclusion; no texts of the Monday feria are spliced into it merely because the feast moved.
\item[Boundary] If another transferred I-class feast is queued, calculate positions and the order of impediment under rubrics 97--98; do not choose the more popular feast.
\item[Rubric locus] 14--15, 23b, 91, 95--99; Mass rubrics 431--432, 475--501.
\end{casefile}

\begin{casefile}{3}{All Souls on Sunday}
\item[Calendar facts] November 2 occurs on a Sunday.  All Souls is a I-class liturgical day but its own rule yields to every Sunday.
\item[Ranks and rows] All Souls is position 8; the occurring ordinary Sunday is position 15.  Rubric 16b is an express exception to a bare comparison of position numbers.
\item[Impeded item] All Souls transfers to Monday, which is treated as its proper seat.  The Sunday is not displaced or commemorated on Monday.
\item[Mass category] Sunday has its Mass of the Office.  Monday has the All Souls Requiem Office and its corresponding Masses, not a private daily Requiem.
\item[Assembly] Assemble the Sunday normally.  On Monday a priest saying one Mass uses the first formulary; one saying two uses the first and second; the three-formulary order and conventual or sung assignment follow rubrics 403--404.  Omit Gloria, Creed, Office commemorations, and blessing; use the Requiem Preface and Requiem variants.
\item[Boundary] The special Monday seat is transfer of the liturgical day, not an external solemnity or an authorization to reproduce All Souls on another convenient date.
\item[Rubric locus] 16b, 91 position 8, 96b, 390--404, 432d, 476f, 499, 507c--510.
\end{casefile}

\begin{casefile}{4}{A church's I-class titular on an ordinary Sunday}
\item[Calendar facts] An approved proper calendar assigns a non-Dominical saint as I-class title of this church, and its date falls on a II-class Sunday.
\item[Ranks and rows] The proper I-class title is position 12; the Sunday is position 15.
\item[Impeded item] The Sunday is privilegedly commemorated in every Mass.  It does not move.
\item[Mass category] The titular feast is the Mass of the Office in this church.  The same rank must not be projected onto other churches or the diocese without their own proper law.
\item[Assembly] Use every proper text of the title and complete missing items from the directed Common.  Use the feast's color, Gloria, Creed, and governing Preface.  Insert the Sunday prayer triad as the sole privileged commemoration.  Do not import the Sunday readings, chants, or Preface.
\item[Boundary] “Titular” and “I class” must be established from the church's approved calendar and title.  A dedication plaque, patronal custom, or parish bulletin is not by itself the calendar source.
\item[Rubric locus] 44, 45b, precedence table 91 positions 12 and 15; 109a, 111a, 113; 305, 431--432, 475--505.
\end{casefile}

\begin{casefile}{5}{A II-class feast of the Lord on a II-class Sunday}
\item[Calendar facts] A universal fixed II-class feast of the Lord, such as the Transfiguration, occurs on an ordinary Sunday.
\item[Ranks and rows] Feasts of the Lord of II class stand at position 14; the Sunday stands at position 15.
\item[Impeded item] The feast replaces the Sunday with all rights and there is no Sunday commemoration.
\item[Mass category] The feast's Mass agrees with the occurring Office; it is not a festive or votive substitution.
\item[Assembly] Use the feast's white proper, completing only explicit Common references.  Say Gloria and Creed.  Determine its Preface by the printed proper and rubrics 482--499.  No Sunday prayer, reading, chant, or Last Gospel is added simply because the civil day is Sunday.
\item[Boundary] The exception concerns feasts of the Lord.  It cannot be transferred to a II-class saint or to a devotion whose subject is the Lord.
\item[Rubric locus] 16a; precedence table 91 positions 14--15; 431--432, 475--499.
\end{casefile}

\begin{casefile}{6}{A II-class saint on a II-class Sunday}
\item[Calendar facts] A universal II-class saint occurs on an ordinary II-class Sunday.
\item[Ranks and rows] The Sunday is position 15; the universal saint is position 16.
\item[Impeded item] The saint is an ordinary commemoration.  It is made at a Low Mass and the conventual Mass, but omitted from a nonconventual sung Mass.  It is not accidentally transferred.
\item[Mass category] The Mass is the Sunday Mass of the Office.
\item[Assembly] Use the full Sunday temporal formulary, seasonal color, Gloria, Creed, seasonal or assigned Preface, and normal conclusion.  At each Low or conventual Mass insert the saint's Collect, Secret, and Postcommunion as the single admitted commemoration; at the nonconventual sung Mass use only the Sunday prayer triad.
\item[Boundary] The two differing prayer sets are not an inconsistency: Mass form controls admission of an ordinary commemoration.
\item[Rubric locus] 16, precedence table 91 positions 15--16; 95, 108, 111b, 113--114; 431--432, 475--505.
\end{casefile}

\begin{casefile}{7}{A III-class saint in early Advent}
\item[Calendar facts] Before December 17, a universal III-class saint and an ordinary Advent feria occur together.
\item[Ranks and rows] The universal III-class feast is position 24; the Advent feria through December 16 is position 25.
\item[Impeded item] The Advent feria is privilegedly commemorated in every Mass.
\item[Mass category] Celebrate the saint's Mass of the Office, not an Advent ferial Mass with a festive commemoration.
\item[Assembly] Use the saint's proper and directed Common in its assigned color.  Gloria follows the Office; there is no Creed merely from III class.  Insert the Advent feria's Collect, Secret, and Postcommunion first.  Choose the Preface from the selected Mass and seasonal rules, never from the commemoration.
\item[Boundary] From December 17 through 23 the Advent feria is II class at position 18 and the result reverses.
\item[Rubric locus] 25b, precedence table 91 positions 18, 24--25; 108--109e, 111d, 113; 270, 431--432, 475--505.
\end{casefile}

\begin{casefile}{8}{A III-class saint on an ordinary Lenten feria}
\item[Calendar facts] From Thursday after Ash Wednesday through Saturday before Palm Sunday, excluding Ember Days, a III-class saint meets a III-class Lenten feria.
\item[Ranks and rows] The Lenten feria is position 22; a particular or universal III-class saint is position 23 or 24.
\item[Impeded item] The saint is ordinarily commemorated only at the conventual and Low Masses; it is absent from a nonconventual sung Mass.
\item[Mass category] The Mass is the Lenten feria's Mass of the Office.
\item[Assembly] Use the violet ferial temporal proper, omit Gloria and Creed, use the Lenten Preface, and append the prayer over the people after all Postcommunions.  Where the saint is admitted, add its Collect, Secret, and Postcommunion after the ferial prayers.  The saint supplies no color, Gospel, or Preface.
\item[Boundary] An Ember day, Ash Wednesday, or Holy Week feria has a different class and position and must be recalculated rather than treated as this case.
\item[Rubric locus] 25a, precedence table 91 positions 7, 18, 22--24; 108, 109e, 111d; 432, 476, 486, 505--506.
\end{casefile}

\begin{casefile}{9}{An Ember day against universal and proper II-class feasts}
\item[Calendar facts] A September Ember day occurs with a II-class feast.  Compare two distinct hypotheses: the feast is universal; or the feast is proper to the place.
\item[Ranks and rows] The universal II-class feast is position 16, the Ember day position 18, and the proper II-class feast position 19.
\item[Impeded item] In the universal-feast hypothesis, the feast wins and the September Ember day is privilegedly commemorated in every Mass.  In the proper-feast hypothesis, the Ember day wins; the accidentally impeded II-class feast does not transfer and is an ordinary commemoration where Mass form and the one-commemoration limit admit it.
\item[Mass category] In each hypothesis use the Mass of the winning Office.
\item[Assembly] Universal winner: feast proper/Common, its Gloria rule, its Creed rule, and its Preface, with the Ember commemoration first; on an Ember Wednesday or Saturday its Collect is the feria's first prayer, the one said at Lauds, under rubric 443.  Ember winner: the violet Ember formulary and ferial Preface, retaining additional lessons only where rubrics 467--468 appoint them; add the proper feast's prayer triad only at Low or conventual Mass.  At a nonconventual sung Mass omit that ordinary commemoration.
\item[Boundary] “Both are II class” does not answer the case; universal and proper provenance determines their rows.
\item[Rubric locus] 24b; precedence table 91 positions 16, 18--19; 95, 108--109d, 111c, 113; 432, 443, 467--468, 475--505.
\end{casefile}

\begin{casefile}{10}{A particular III-class saint against a universal III-class saint}
\item[Calendar facts] An approved diocesan calendar places a genuinely particular III-class saint on the date of a different universal III-class saint.
\item[Ranks and rows] Particular III-class feasts are position 23; universal III-class feasts are position 24.
\item[Impeded item] The universal saint becomes an ordinary commemoration where admitted; neither feast acquires accidental transfer.
\item[Mass category] The particular saint's Mass is the Mass of the Office within that calendar only.
\item[Assembly] Use the particular proper, or the exact Common to which it refers, with its color and Gloria rule; ordinarily omit Creed.  At Low and conventual Masses add the universal saint's prayer triad, subject to the two-commemoration ceiling.  A nonconventual sung Mass omits the ordinary commemoration.
\item[Boundary] Without the promulgated diocesan calendar or supplement the asserted proper rank is unproved, so the case has no determinate local answer.  Never infer position 23 merely from local affection.
\item[Rubric locus] 42--46; precedence table 91 positions 23--24; 95, 108, 111d, 113; 305, 431--432, 475--505.
\end{casefile}

\begin{casefile}{11}{A II- or III-class vigil is impeded}
\item[Calendar facts] A vigil of II or III class meets another liturgical day.
\item[Ranks and rows] II-class vigils are position 21; the III-class Vigil of Saint Lawrence is position 26.
\item[Impeded item] On any Sunday or I-class feast, the vigil is wholly omitted.  It is also omitted if its feast is transferred or reduced to commemoration.  Against another superior non-Sunday day, an impeded vigil is commemorated according to the ordinary-commemoration rules.
\item[Mass category] Use the winning day's Mass of the Office.  A vigil is not an “anticipated feast Mass.”
\item[Assembly] When omission applies, insert nothing from the vigil.  When commemoration applies, add only its prayer triad at the conventual and Low Masses and only within that day's numerical limit; a nonconventual sung Mass omits it.  The winner controls color, Gloria, Creed, readings, chants, Preface, and Last Gospel.
\item[Boundary] The I-class Christmas and Pentecost vigils have their own superior rows and are not examples of this omission rule.
\item[Rubric locus] 28--34; precedence table 91 positions 5, 9, 21, 26; 108, 111--114.
\end{casefile}

\begin{casefile}{12}{The special days inside the Christmas octave}
\item[Calendar facts] First take a weekday from December 26 through 31; then take a Sunday in that range.
\item[Ranks and rows] The octave days are position 17.  Saint Stephen, Saint John, and the Holy Innocents are universal II-class feasts at position 16; the Sunday is position 15.  Rubrics 68--70 add special rules.
\item[Impeded item] On their weekdays the three named feasts are celebrated.  On Sunday, the Sunday is always celebrated and the occurring II-class feast is ordinarily commemorated.  Saint Thomas Becket and Saint Sylvester are commemorated on their assigned weekdays.  Under rubric 68f's special rule for all particular feasts, only a particular I-class feast honoring a saint already inscribed on that date is admitted; every other particular feast, including one of II or III class, transfers after the octave.
\item[Mass category] In each branch celebrate the Mass corresponding to the Office selected by the octave rules.
\item[Assembly] Retain the day's or feast's proper formulary and appointed color---red for Saint Stephen and the Holy Innocents, white for Saint John, the Sunday, and a bare octave day. Apply the admitted prayer triad, and use the Christmas Preface and proper Christmas-octave \latin{Communicantes} where rubrics 484 and 501 require them, even when a saint's Mass is celebrated.
\item[Boundary] “The octave is II class” is insufficient; the date-specific structure of rubrics 68--70 controls.
\item[Rubric locus] 64--70; precedence table 91 positions 15--17; 109c, 111; 484, 501.
\end{casefile}

\begin{casefile}{13}{A proper I-class feast inside the Easter octave}
\item[Calendar facts] The I-class title or dedication of a church falls on a weekday of the Easter octave.
\item[Ranks and rows] The octave day is position 10; the proper I-class feast is position 12.
\item[Impeded item] The proper feast transfers to the nearest following day that is not I or II class.  The octave day is celebrated without borrowing the feast's texts.
\item[Mass category] First celebrate the temporal octave Mass of the Office; later celebrate the transferred feast as a I-class Mass of the Office.
\item[Assembly] The octave Mass uses its white temporal proper, Gloria, Creed, Paschal dismissal Alleluias, Easter Preface, and proper octave Canon forms.  The feast on its transferred day uses its own proper/Common and color, Gloria and Creed.  If still inside an octave in another hypothesis, rubric 501's Canon overlay would have to be applied independently.
\item[Boundary] If several I-class feasts wait, rubrics 97--98 govern their order.  The local title or dedication and rank must be sourced.
\item[Rubric locus] 66, precedence table 91 positions 10 and 12; 95--99; 431--432, 475--476, 489, 501, 507b.
\end{casefile}

\begin{casefile}{14}{Accidental transfer versus perpetual reposition}
\item[Calendar facts] A proper II-class feast is impeded once by a movable superior day; a second proper II-class feast is fixed so that the same superior day impedes it every year throughout the relevant calendar.
\item[Ranks and rows] Classify each conflict from the precedence table; class equality alone is not the test.
\item[Impeded item] The accidentally impeded II-class feast has no general right of transfer and is commemorated or omitted as the other rules direct.  The perpetually impeded II-class feast has a right of stable reposition to the nearest following day not I or II class; that new day becomes its proper day.
\item[Mass category] On the impeding day use its Mass of the Office.  Only the feast lawfully reposed in the approved calendar has a Mass of the Office on the new date.
\item[Assembly] The accidental case contributes at most an admitted prayer triad.  The reposed case uses the feast's complete proper/Common, class, color, Gloria and Creed rules, and Preface on its stable new date.
\item[Boundary] The manual cannot choose a new date for a real diocese or church.  The competent calendar must record the reposition and resolve collisions with already existing entries.
\item[Rubric locus] 92--102, especially 95 and 100--102; 106--114; 269--305.
\end{casefile}

\begin{casefile}{15}{Too many commemorations on a I-class day}
\item[Calendar facts] A I-class feast wins, while a Sunday, an ordinary saint, and an otherwise requested additional prayer are proposed.
\item[Ranks and rows] The winner's exact position is established; the Sunday is a privileged commemoration, while the saint is ordinary.
\item[Impeded item] Only one privileged commemoration is admitted.  The Sunday prayer triad is retained; the ordinary saint and additional prayer are omitted.
\item[Mass category] The Mass is the I-class Mass of the Office, not a composite of the three observances.
\item[Assembly] Use the winner's entire formulary, Gloria, Creed, color, and Preface.  Add the Sunday Collect, Secret, and Postcommunion first and alone.  Nothing else changes the readings, chants, Preface, Canon, or Last Gospel.
\item[Boundary] The same one-privileged-only limit applies to a I-class votive and a nonconventual sung Mass, but the reason for the limit must still be recorded.
\item[Rubric locus] 106--114, especially 109a, 111a, 113--114; Mass limit 433--446.
\end{casefile}

\begin{casefile}{16}{A II-class Sunday, a II-class saint, and the Major Litanies}
\item[Calendar facts] An ordinary II-class Sunday coincides with a II-class saint and with the Major Litanies, in a setting where the Sunday Mass rather than the Rogation votive is the admitted Mass.
\item[Ranks and rows] Sunday position 15 defeats the saint at position 16 or 19.  The Major Litanies create a privileged commemoration in Mass.
\item[Impeded item] A II-class Sunday ordinarily admits one II-class feast, but rubric 111b omits that feast when a privileged commemoration must be made.  Thus the Litanies prayer displaces the saint's commemoration.
\item[Mass category] The Sunday remains the Mass of the Office.  If a procession makes the Rogation Mass applicable, recalculate under rubrics 346--347 rather than reusing this answer.
\item[Assembly] Use the full Sunday proper, Gloria, Creed, and governing Preface.  Add only the Major-Litanies prayer triad.  Omit the saint completely from this Mass.
\item[Boundary] “One commemoration allowed” does not mean the celebrant may choose which one; privileged status decides.
\item[Rubric locus] 109f, 111b, 113--114; 346--347; precedence table 91 positions 15--16 and 19.
\end{casefile}

\begin{casefile}{17}{Competing commemorations on another II-class day}
\item[Calendar facts] A universal II-class feast prevails over a September Ember day, while another proper III-class saint is also entered in the sourced calendar.
\item[Ranks and rows] The winner is position 16; the Ember day is position 18 and privileged; the proper III-class saint is position 23 and ordinary.
\item[Impeded item] A non-Sunday II-class day admits only one commemoration, privileged or ordinary.  The Ember feria is kept; the saint is omitted.
\item[Mass category] Celebrate the universal feast's Mass of the Office.
\item[Assembly] Use the feast's proper/Common, color, Gloria and Creed rules, and Preface.  Add the Ember Collect, Secret, and Postcommunion as the one commemoration.  The Ember readings and ferial Preface do not enter the feast Mass.
\item[Boundary] If the winning feast itself is merely proper II class, the Ember day would precede it and the entire calculation would reverse.
\item[Rubric locus] Precedence table 91 positions 16, 18, 23; 109d, 111c, 113--114; 433--446, 480, 505.
\end{casefile}

\begin{casefile}{18}{Two commemorations on a III-class day and a duplicate exclusion}
\item[Calendar facts] A particular III-class saint wins in early Advent over both the Advent feria and a different universal III-class saint.  A free prayer invoking the winning saint is also proposed.
\item[Ranks and rows] The particular feast is position 23, the universal feast position 24, and the Advent feria position 25.
\item[Impeded item] The Advent feria is the first, privileged commemoration; the universal saint is the second, ordinary commemoration where Mass form admits it.  The extra prayer invoking the winning saint is excluded by identity before numerical counting.
\item[Mass category] The particular saint's Mass is the Mass of the Office in the sourced proper calendar.
\item[Assembly] At a Low or conventual Mass use the winner's proper/Common, then the Advent and universal-saint prayer triads in that order, never exceeding three total prayers.  At a nonconventual sung Mass ordinary commemorations are barred, so only the Advent triad remains.  The winner alone supplies chants, readings, color, and Preface.
\item[Boundary] The same-Person exclusion is not a discretionary way to make room; it prevents duplicate liturgical invocation even when the numerical ceiling has space.
\item[Rubric locus] 108--114, especially 109e, 111a/d, 112d, 113; precedence table 91 positions 23--25; 433--446.
\end{casefile}
